% Generated draft: requires mathematical and source review.
\documentclass[11pt]{article}
\usepackage[margin=25mm]{geometry}
\usepackage{fontspec}
\setmainfont{Latin Modern Roman}
\usepackage{amsmath,amssymb,mathtools,unicode-math}
\setmathfont{Latin Modern Math}
\usepackage{enumitem,needspace,xurl,hyperref}
\setlength{\parindent}{0pt}
\setlength{\parskip}{5pt}
\setlength{\emergencystretch}{3em}
\newcommand{\N}{\mathbb{N}}
\newcommand{\Z}{\mathbb{Z}}
\newcommand{\Q}{\mathbb{Q}}
\newcommand{\R}{\mathbb{R}}
\newcommand{\C}{\mathbb{C}}
\newcommand{\F}{\mathbb{F}}
\newcommand{\A}{\mathbb{A}}
\newcommand{\PP}{\mathbb{P}}
\newcommand{\E}{\mathbb{E}}
\newcommand{\eps}{\varepsilon}
\newcommand{\dd}{\,\mathrm{d}}
\begin{document}
\section*{UnsolvedMath ID: 30000909}
\textbf{Source problem code:} catalogue entry
\subsection{Definitions and mathematical statement}
For \(n\geq1\), let \(F_n\) be the free group with its ordered free generators, and identify an \(n\)-marked group with a quotient \(F_n/N\) for a normal subgroup \(N\triangleleft F_n\). The space \(\mathcal G_n\) of such normal subgroups is a closed subspace of \(\{0,1\}^{F_n}\), equivalently carrying the topology of agreement on larger finite marked Cayley balls. Use the customary Polish coding
\[
\mathcal G=\bigsqcup_{n\geq1}\mathcal G_n
\]
for finitely generated marked groups. Repeated or redundant generators are allowed by this quotient coding.

Equip each group with the word metric of its finite marking. Two groups \(G,H\) are quasi-isometric if there exist a map \(q:G\to H\) and constants \(\lambda\geq1\), \(C\geq0\) such that
\[
\lambda^{-1}d_G(x,x')-C\leq d_H(q(x),q(x'))
\leq\lambda d_G(x,x')+C
\quad(x,x'\in G),
\]
and every point of \(H\) lies within distance \(C\) of \(q(G)\). Write \(E_{\mathrm{QI}}\) for this equivalence relation on \(\mathcal G\); its definition concerns the groups with word metrics, rather than an isomorphism respecting the markings.

An equivalence relation \(E\subseteq Y\times Y\) on a Polish space is \(K_\sigma\) if it is a countable union of compact subsets of \(Y\times Y\). A Borel reduction of \(E\) to \(E_{\mathrm{QI}}\) is a Borel map \(r:Y\to\mathcal G\) such that
\[
y\mathrel E y'\quad\Longleftrightarrow\quad
r(y)\mathrel{E_{\mathrm{QI}}}r(y')
\qquad(y,y'\in Y).
\]
The conjecture asks whether \(E_{\mathrm{QI}}\) is a universal \(K_\sigma\) Borel equivalence relation: it belongs to that class, and every \(K_\sigma\) Borel equivalence relation on a Polish space has such a reduction to it. The reduction may depend on \(E\), and it need not be continuous or injective. The supplied assessment reports the \(K_\sigma\) upper bound as known; the target left open there is universality among finitely generated groups, not merely among bounded-degree graphs.
\subsection{Short English statement}
Is quasi-isometry of finitely generated marked groups as complicated under Borel reduction as every \(K_\sigma\) Borel equivalence relation? The source reports the upper bound but conjectures the matching universality.
\subsection{Sources}
\begin{itemize}
\item Local UnsolvedMath registry, permanent ID \(30000909\), source code catalogue entry, ``Universality of Quasi-Isometry among Marked Groups.'' No individual proposer is supplied.
\item \emph{Set Theory} (2008), Oberwolfach Reports, DOI 10.4171/owr/2008/02, pp. 79--118: \url{https://doi.org/10.4171/owr/2008/02}.
\item Supplied reference, ``Thomas'': \url{https://sites.math.rutgers.edu/~sthomas/quasi2.pdf}.
\item Supplied literature triage, checked 2026-08-21: status ``open.'' It attributes to Thomas the \(K_\sigma\) upper bound and a universal relation for connected four-regular graphs, while calling quasi-isometry universality for finitely generated groups conjectural. This dated assessment derives from supplied OWR curation; time-sensitive status and cited claims require independent verification.
\end{itemize}
\end{document}
